% Owner: Yashpal (Member 1)
\subsection{Reusable cache}
One module, \sig{cache}, serves as the I-cache, the D-cache and the unified cache
(Fig.~\ref{fig:cache}). It is direct-mapped, with \sig{NUM\_LINES} lines of \sig{LINE\_WORDS}
32-bit words (both powers of two), one valid bit and one tag per line. The fields are computed
with shifts and masks, so a one-line cache or a one-word line also works. All results use
4-word lines.
\begin{itemize}
  \item \textbf{Read hit:} completes in the same cycle through a combinational array read.
  \item \textbf{Read miss:} requests the whole line, writes it into the array in the cycle the
        memory answers, and completes as a hit in the next cycle. The miss penalty is
        therefore exactly $L$ cycles. A read miss first waits for the write buffer to drain,
        so the line it fetches already contains any buffered store (no read-after-write
        hazard).
  \item \textbf{Writes:} write-through and no-write-allocate. A store hit also updates the
        word in the line. A store miss leaves the cache unchanged.
  \item \textbf{Write buffer} (\sig{WBUF}\,=\,1): one entry. A store that finds the buffer
        empty is offered to memory at once. If memory does not finish it in that cycle, the
        buffer captures it and keeps the same request asserted, so the latency count simply
        continues while the core goes on. A store that finds the buffer busy waits until it
        empties. With $L=1$ the buffer is therefore never used, and a store is never slower
        than without the buffer. Without the buffer a store waits for memory.
  \item \textbf{Uncached mode} (\sig{ENABLE}\,=\,0): every read goes to memory and its word is
        passed straight back. With $L=1$ this is the Phase~I combinational memory.
\end{itemize}
The cache also outputs event pulses for refills and write-buffer stalls, which the performance
counters use.

\begin{figure}[t]
  \centering
  \resizebox{0.92\columnwidth}{!}{% Organisation of the reusable direct-mapped cache (rtl/mem/cache.v)
\begin{tikzpicture}[
    font=\footnotesize,
    fld/.style={draw, thick, minimum height=6mm, align=center, inner sep=2pt},
    blk/.style={draw, thick, rounded corners=2pt, minimum height=7mm, align=center, fill=white},
    arr/.style={-{Stealth[length=2mm]}, thick}]
  % address fields
  \node[fld, minimum width=22mm, fill=orange!20] (tag) at (0,0) {tag};
  \node[fld, minimum width=14mm, fill=green!20, right=0mm of tag] (idx) {index};
  \node[fld, minimum width=12mm, fill=blue!15, right=0mm of idx] (off) {offset};
  \node[fld, minimum width=7mm, right=0mm of off] (bb) {00};
  \node[anchor=south, font=\scriptsize] at (1.8,0.35) {byte address (low \sig{ABITS} bits)};
  % arrays
  \node[blk, minimum width=9mm, minimum height=16mm] (v) at (-0.2,-2.0) {V};
  \node[blk, minimum width=14mm, minimum height=16mm, right=1mm of v] (t) {tags};
  \node[blk, minimum width=30mm, minimum height=16mm, right=1mm of t] (d) {lines\\$N \times W$ words};
  \draw[arr] (idx.south) -- ++(0,-0.45) -| (t.north);
  \draw[arr] (idx.south) ++(0,-0.45) -| (d.north);
  \draw[arr] (idx.south) ++(0,-0.45) -| (v.north);
  \node[blk, minimum width=10mm] (cmp) at (0.5,-3.65) {$=$};
  \draw[arr] (t.south) -- (cmp.north -| t.south) -- (cmp.east);
  \draw[arr] (tag.west) -- ++(-0.35,0) |- (cmp.west);
  \node[blk, minimum width=8mm] (and) at (-0.65,-4.5) {\&};
  \draw[arr] (v.south) -- (and.north);
  \draw[arr] (cmp.south) |- (and.east);
  \node[anchor=north, font=\scriptsize\ttfamily] at (and.south) {hit};
  \node[blk, minimum width=12mm] (mux) at (3.1,-3.65) {word\\select};
  \draw[arr] (d.south -| mux.north) -- (mux.north);
  \draw[arr] (off.south) -- ++(0,-0.25) -| (5.25,-3.65) -- (mux.east);
  \node[anchor=north, font=\scriptsize\ttfamily] at (mux.south) {rdata};
  % write buffer and memory side
  \node[blk, minimum width=18mm, fill=blue!8] (wb) at (7.4,-0.55) {write buffer\\(1 entry)};
  \node[blk, minimum width=18mm, fill=gray!12] (mp) at (7.4,-2.0) {memory port\\req / ready};
  \draw[arr] (wb) -- node[right, font=\scriptsize, align=left] {drains\\first} (mp);
  \draw[arr] (mp.west) -- node[above, font=\scriptsize] {refill} (d.east |- mp.west);
\end{tikzpicture}
}
  \caption{The reusable direct-mapped cache: address fields, valid/tag/line arrays, hit logic,
  word select, and the one-entry write buffer in front of the memory port.}
  \label{fig:cache}
\end{figure}

\subsection{Shared lower-level memory}
\sig{main\_mem} is a 512-word single-port memory. A requester raises \sig{req} with
\sig{we}, \sig{addr} and \sig{wdata} and holds them until \sig{ready}, which rises in the $L$-th
cycle of the request. A read returns the aligned 4-word line that contains \sig{addr}, so a
refill is one transfer. A write stores one word. With $L=1$, \sig{ready} rises in the same cycle
as \sig{req}, which is exactly the ideal memory of Phase~I. In synthesis the memory is a black
box that stands for block RAM or off-chip memory.
